\section{The Pair Problem}

The previous chapter left the grammar with one concrete reader for the
positive invariant. The invariant was no longer a floating name. It had an
anchored operator, a scale, a kernel of labels, and a rule for deciding when a
sensor class could carry a reading. The present chapter asks what that
reader is forced to do when the first charged difference appears. The answer
cannot be imported as a particle story. It must be read from the grammar that
has already been earned.

The pair problem begins at the point where stationarity is too poor to
finish the record. The first variation has vanished. The null drift has been
pinned. The second variation is the surviving activity. But a quadratic
activity can be read in more than one face. It can be read as magnitude. It
can be read as strain. It can also be read as an oriented charge when the
baseline against which the activity is compared has been fixed. Pair
production is the grammar's name for the moment when that oriented reading
cannot be avoided.

This is why the chapter does not begin with two preexisting objects. A pair
is not placed into the ledger and then assigned signs. The signs are forced
by the way the ledger reads a coupled disturbance. If the coupled residue
telescopes away, no charged distinction remains. If the residue refuses to
cancel under the permitted action, the grammar must discriminate it. The
discrimination is signed because it is read relative to a baseline. The
baseline is not decorative. It decides which orientation is native and which
orientation appears only after a tilt.

The chapter moves through seven obligations. First, the electron emerges as
the flat-baseline reading of the discriminating action, and the charge is
\(-1\). Second, the positron appears only when the same pair is read against
a tilted baseline, where the sign is forced to \(+1\). Third, the chapter
isolates the two identity grants on which the physics rests: the local
quotient that lets selection stand after collapse, and the motion grant that
lets stationarity answer to the world. Fourth, boundary geometry generates
the gauge by producing the action that discriminates the pair residue. Fifth,
the geometric and gauge readings split; a finite carrier may read one side or
the other, but it may not form an interior direct sum that hides a mixed
scalar. Sixth, the electromagnetic reading becomes the Cauchy limit of the
geometric reading as the mesh residue is forced to zero. Seventh, finite
Cohen forcing names the disciplined ledger of bounded spline conditions by
which the continuum reading is compelled without being assumed.

Examples will appear, but only as lanterns. A coincidence detector may
illustrate the cancellation of an electron and a positron. A loop holonomy
may illustrate why an open path carries no charge while a closed path carries
a sign. A compact disc reader may illustrate a threshold gauge, and the
Cantor--G\"odel--Cohen warning may illustrate why completion cannot be
accepted without recoverability. None of these examples supplies the spine.
The spine is the grammatical obligation: what the finite reader permits,
what it forbids, what it identifies, what it quotients, what it carries, and
what it is forced to read at the boundary.

The title therefore names one process, not two. Pair production and geometry
generating gauge are the same turn seen from opposite sides. From the charge
side, the second variation discriminates a pair and forces the electron
reading. From the boundary side, geometry is not a container but the
stabilized grammar of separations; when that boundary grammar is asked to
carry the discriminating residue, it generates the gauge that reads the pair.
The chapter's task is to keep those two faces together without collapsing
them into an illicit direct sum.

